% ============================================================
%  Dispensa 03: Curva EIOPA Risk-Free — Metodo Smith-Wilson
%               Approccio variazionale (struttura stile PCA)
%  Corso di Calcolo Numerico — Laboratorio
%  Università degli Studi di Verona — A.A. 2026–2027
% ============================================================
\documentclass[11pt,a4paper]{article}

\usepackage[T1]{fontenc}
\usepackage[utf8]{inputenc}
\usepackage[italian]{babel}
\usepackage{lmodern}
\usepackage{amsmath,amssymb,amsthm,mathtools}
\usepackage{geometry}
\usepackage{booktabs}
\usepackage{array}
\usepackage{enumitem}
\usepackage{graphicx}
\usepackage[table]{xcolor}
\usepackage{hyperref}
\usepackage{float}
\usepackage{fancyvrb}
\usepackage{tcolorbox}
\tcbuselibrary{breakable}
\usepackage{titlesec}

% le figure sono prodotte da R/eiopa_rfr_sw_variazionale.R in output/eiopa_rfr_sw_var/
\graphicspath{{}}

\geometry{top=2.5cm,bottom=2.5cm,left=2.8cm,right=2.8cm}
\setlist[itemize]{topsep=4pt,itemsep=2pt,parsep=0pt}
\setlist[enumerate]{topsep=4pt,itemsep=2pt,parsep=0pt}
\hypersetup{colorlinks=true,linkcolor=blue!55!black,urlcolor=blue!55!black}

\titleformat{\section}{\large\bfseries}{\thesection.}{0.55em}{}[{\vspace{1pt}\titlerule[0.4pt]}]
\titleformat{\subsection}{\normalsize\bfseries}{\thesubsection.}{0.55em}{}
\titleformat{\subsubsection}{\small\bfseries\itshape}{\thesubsubsection.}{0.5em}{}

%--- ambienti teorema (solo enunciati, niente dimostrazioni)
\theoremstyle{plain}
\newtheorem{teorema}{Teorema}[section]
\newtheorem{proposizione}[teorema]{Proposizione}
\theoremstyle{definition}
\newtheorem{definizione}[teorema]{Definizione}
\newtheorem{esempio}[teorema]{Esempio}
\theoremstyle{remark}
\newtheorem*{osservazione}{Osservazione}
\newtheorem*{attenzione}{Nota numerica}

%--- macro acronimi (funzionano dentro e fuori math mode)
\newcommand{\UFR}{\ensuremath{\mathrm{UFR}}}
\newcommand{\LLP}{\ensuremath{\mathrm{LLP}}}
\newcommand{\CP}{\ensuremath{\mathrm{CP}}}
\newcommand{\CRA}{\ensuremath{\mathrm{CRA}}}
\newcommand{\Qmat}{\mathbf{Q}}
\newcommand{\Hmat}{\mathbf{H}}
\newcommand{\transp}{^{\!\top}}

%--- box colorati
\tcbset{
  mybase/.style={
    colback=gray!5, colframe=gray!40,
    boxrule=0.4pt, arc=2pt,
    left=7pt, right=7pt, top=5pt, bottom=5pt,
    fonttitle=\small\bfseries
  }
}
\newtcolorbox{keybox}[1][Concetto chiave]{mybase,
  colframe=blue!40!black, colback=blue!4, title={#1}}
\newtcolorbox{algobox}[1][Algoritmo]{mybase,
  colframe=teal!50!black, colback=teal!4, title={#1}}
\newtcolorbox{warningbox}{mybase,
  colframe=orange!70!black, colback=orange!5}
\newtcolorbox{codebox}[1][Codice R]{mybase,
  colframe=green!40!black, colback=green!3, title={#1}}

% ============================================================
\title{%
  \textbf{Costruzione della Curva EIOPA Risk-Free}\\[6pt]
  \large Il metodo Smith--Wilson: un approccio variazionale\\[4pt]
}
\author{Laboratorio di Calcolo Numerico --- Laurea Triennale in Matematica Applicata\\
  Universit\`a degli Studi di Verona\\[2pt]
  Anno Accademico 2026--2027}
\date{}

\begin{document}
\maketitle
\tableofcontents
\newpage

% ============================================================
\section{Motivazioni economiche}
% ============================================================

\subsection{Il problema: valutare oggi impegni decennali}

Una compagnia di assicurazione vita o un fondo pensione assume \textbf{obblighi di pagamento
che si estendono per decenni}: una rendita vitalizia o una polizza possono generare uscite
fino a 50--60 anni nel futuro.

Il quadro regolamentare europeo \textbf{Solvency~II} chiede che ogni compagnia dimostri,
\emph{oggi}, di avere risorse sufficienti per far fronte a tutti questi impegni futuri ---
\`e il principio di \textbf{solvibilit\`a}. Per farlo occorre rispondere alla domanda:
\emph{quanto valgono oggi quei pagamenti futuri?}

La risposta dipende dal \textbf{valore temporale del denaro}: 100\,\texteuro{} ricevuti tra
30 anni valgono \emph{meno} di 100\,\texteuro{} ricevuti oggi, perch\'e nel frattempo il
denaro pu\`o essere investito e generare interessi. Quanto meno? Dipende dal
\textbf{tasso di interesse} applicabile su quell'orizzonte temporale.

Poich\'e le passivit\`a hanno scadenze diverse, serve conoscere un tasso di interesse per
ogni orizzonte: 1 anno, 5 anni, 20 anni, 50 anni\ldots\ L'insieme di questi tassi al
variare della scadenza $T$ \`e la \textbf{struttura a termine dei tassi di interesse}, o
pi\`u brevemente \textbf{curva dei tassi}.

Solvency~II impone che tutte le imprese europee usino la \textbf{stessa} curva, calcolata e
pubblicata mensilmente da \textbf{EIOPA} (European Insurance and Occupational Pensions
Authority). La metodologia ufficiale \`e descritta nel documento tecnico
EIOPA-BoS-25-599~\cite{eiopa2025}.

\subsection{Il problema dei dati: il mercato finisce a 20 anni}

Il mercato finanziario fornisce quotazioni di tassi di interesse per molte scadenze tramite
contratti standardizzati. Per l'euro questi contratti sono \textbf{attivamente scambiati e
affidabili} fino a circa 20 anni di scadenza: questo orizzonte \`e detto
\textbf{Last Liquid Point} ($\LLP = 20$ anni). Oltre tale soglia le quotazioni diventano
sporadiche, poco rappresentative o del tutto assenti.

Il problema \`e evidente: le passivit\`a assicurative si estendono ben oltre i 20 anni,
ma i dati di mercato si esauriscono proprio l\`i. EIOPA deve comunque fornire una stima
del tasso di interesse anche per le scadenze lunghe, fino al \textbf{Convergence Point}
($\CP = 60$ anni).

\begin{keybox}[Due problemi in uno]
Costruire la curva EIOPA significa risolvere due problemi distinti:
\begin{enumerate}
  \item \textbf{Ricostruire} in modo regolare i tassi alle scadenze dove esistono
        quotazioni di mercato ($T \le \LLP = 20$ anni);
  \item \textbf{Estendere} la curva oltre il $\LLP$ fino al $\CP = 60$ anni in modo
        stabile, coerente e non arbitrario.
\end{enumerate}
\end{keybox}

\subsection{L'ancora di lungo periodo: l'UFR}

Come estendere la curva oltre i 20 anni? L'idea ingenua --- prolungare l'ultimo tasso
osservato ``a piatto'' --- \`e fragile: una piccola variazione nel tasso a 20 anni si
propagherebbe su tutti gli anni da 21 a 60, amplificando le oscillazioni nei bilanci
assicurativi.

La scelta regolatoria (\cite{eiopa2025}, sez.~3.6--3.8) \`e di \textbf{agganciare la curva
a un tasso di equilibrio di lungo periodo}: l'\textbf{Ultimate Forward Rate}
($\UFR = 3{,}30\%$ per l'euro). L'$\UFR$ rappresenta la stima regolamentare del tasso
reale di crescita economica di lungo periodo sommato all'inflazione attesa; cambia molto
lentamente nel tempo ed \`e quindi un'\textbf{ancora stabile}.

La curva viene costruita in modo che il tasso di interesse, superato il $\LLP$, converga
\emph{gradualmente e in modo regolare} verso l'$\UFR$ entro il $\CP = 60$ anni.

\begin{figure}[H]
\centering
\includegraphics[width=0.92\textwidth]{../output/eiopa_rfr_sw_var/fig1_motivazione_estrapolazione.pdf}
\caption{La curva EUR calcolata da EIOPA (dicembre~2025) fino al Convergence Point di 60 anni.
Nella zona ombreggiata ($T > \LLP = 20$ anni) non esistono dati di mercato affidabili: la
curva \`e costruita in modo che il tasso converga al valore di lungo periodo
$\UFR = 3{,}30\%$. \`E questo l'obiettivo del metodo Smith--Wilson.}
\label{fig:motivazione}
\end{figure}

\subsection{Eliminare il rischio di controparte: il CRA}
\label{sec:cra}

I contratti di mercato usati come sorgente dati (swap di tasso di interesse) includono un
piccolo \textbf{premio per il rischio di controparte}: chi riceve i pagamenti sopporta il
rischio che la controparte non onori il contratto. Questo premio non riflette il valore
temporale del denaro in s\'e, e deve essere rimosso.

Per ottenere un tasso \emph{privo di rischio} EIOPA sottrae un \textbf{Credit Risk
Adjustment} ($\CRA = 10$ bps per l'euro, \cite{eiopa2025} sez.~3.5) dai tassi di mercato
\emph{prima} di calibrare la curva. Tutti i tassi di input che useremo nella
Sez.~\ref{sec:ricostruzione} sono gi\`a corretti di questi 10 punti base.

% ============================================================
\section{La curva dei tassi e i tassi forward}
% ============================================================

Tutte le nozioni di tasso --- spot, forward, par --- sono contenute in un unico oggetto:
il \textbf{fattore di sconto} $P(0,t)$. Spot e forward sono semplicemente due modi diversi
di ``leggere'' questa stessa funzione. Costruiamo la sua definizione partendo dalla
capitalizzazione discreta e arrivando alla capitalizzazione continua come caso limite naturale.

% ------------------------------------------------------------
\subsection{Dalla capitalizzazione al fattore di sconto}
% ------------------------------------------------------------

\paragraph{Capitalizzazione annua composta.}
\emph{Capitalizzare} significa trasformare un euro di oggi in pi\`u euro in futuro.
Investendo $1$~\texteuro{} a tasso $r$ per $T$ anni si ottiene il
\textbf{fattore di capitalizzazione}
\[
  C_{\rm ann}(r,T) = (1+r)^T.
\]
Ad esempio con $r=3\%$ e $T=10$ anni: $C = 1.03^{10} \approx 1.344$, cio\`e $1$~\texteuro{}
cresce a $1.34$~\texteuro{}.

\paragraph{Capitalizzazione $m$ volte l'anno.}
Se la capitalizzazione avviene $m$ volte per anno al tasso nominale $r_m$ (rate $r_m/m$
per periodo), il fattore di capitalizzazione diventa
\[
  C_m(r_m,T) = \left(1 + \frac{r_m}{m}\right)^{mT}.
\]
All'aumentare di $m$ la capitalizzazione diventa sempre pi\`u frequente: $m=12$ mensile,
$m=365$ giornaliera, e cos\`i via.

\paragraph{Limite continuo: $m \to \infty$.}
Usando il limite fondamentale $\lim_{m\to\infty}(1+x/m)^m = e^x$ si ottiene il
\textbf{fattore di capitalizzazione continua}:
\begin{equation}
  C_\infty(r,T) = \lim_{m\to\infty}\left(1+\frac{r_m}{m}\right)^{mT} = e^{r\,T}.
  \label{eq:cap-cont}
\end{equation}
La capitalizzazione continua \`e lo standard della matematica finanziaria quantitativa
perch\'e l'esponenziale rende i \emph{logaritmi} dei fattori additivi (invece che i fattori
stessi moltiplicativi), aprendo la strada alla derivata che definisce il tasso forward.

\paragraph{Il fattore di sconto: il problema inverso.}
Finora abbiamo chiesto: \emph{``quanto vale in futuro un euro investito oggi?''} Il problema
\emph{inverso} --- \emph{``quanto vale oggi un euro che riceveremo al tempo $t$?''} --- \`e
il problema dello sconto. Il fattore di sconto \`e semplicemente il reciproco del fattore
di capitalizzazione continua.

\begin{definizione}[Fattore di sconto]
\label{def:sconto}
Il \textbf{fattore di sconto} $P(s,t)$ \`e il valore al tempo $s$ di $1$~\texteuro{}
pagato al tempo $t \ge s$:
\[
  P(s,t) = \frac{1}{C_\infty(r,\,t-s)} = e^{-r(s,t)\,(t-s)}.
\]
Il primo argomento $s$ \`e la \textbf{data di osservazione} (quando valutiamo), il secondo
$t$ \`e la \textbf{data di pagamento} (quando riceviamo l'euro). Nel seguito fissiamo
$s = 0$ (oggi) e scriviamo semplicemente $P(0,t)$.

Propriet\`a immediate:
\begin{itemize}
  \item $P(0,0) = 1$ (un euro pagato subito vale un euro);
  \item $P(0,t)$ strettamente decrescente in $t$ se $r > 0$;
  \item tutta la struttura a termine dei tassi \`e codificata in $P(0,\cdot)$.
\end{itemize}
\end{definizione}

% ------------------------------------------------------------
\subsection{Tasso spot: leggere $P(0,t)$}
% ------------------------------------------------------------

Dalla Definizione~\ref{def:sconto} il fattore di sconto ha la forma $P(0,t)=e^{-r(t)t}$.
Invertendo si ricava il tasso che lo genera.

\begin{definizione}[Tasso spot]
\label{def:spot}
Il \textbf{tasso spot} $r(t)$ \`e il tasso in capitalizzazione continua associato al
fattore di sconto $P(0,t)$:
\begin{equation}
  r(t) = -\frac{\ln P(0,t)}{t}.
  \label{eq:spot-inv}
\end{equation}
\`E la media logaritmica del costo del denaro sull'intero intervallo $[0,t]$.
\end{definizione}

\paragraph{Perch\'e il logaritmo rende tutto pi\`u semplice.}
In capitalizzazione discreta i fattori di sconto si \emph{moltiplicano}: il fattore su
$[0,t_1+t_2]$ \`e il prodotto di quello su $[0,t_1]$ e di quello su $[t_1,t_1+t_2]$,
il che rende le composizioni scomode da derivare. In capitalizzazione continua, invece,
$\ln P(0,t) = -r(t)\cdot t$ \`e \emph{additivo}: i log-fattori si sommano, e questo
\`e il punto di ingresso naturale per la derivata che definisce il forward.

\paragraph{Conversione con la capitalizzazione annua.}
Il tasso annuo composto $r_{\rm ann}(t) = P(0,t)^{-1/t}-1$ \`e legato al tasso continuo da
\[
  r_{\rm ann}(t) = e^{r(t)} - 1, \qquad r(t) = \ln(1 + r_{\rm ann}(t)).
\]
EIOPA pubblica i risultati in capitalizzazione annua (i par swap pagano cedole annue),
ma il metodo Smith--Wilson costruisce $P(0,t)$ internamente usando la capitalizzazione
continua: la conversione $r_{\rm ann} = e^r - 1$ avviene solo all'output finale.

% ------------------------------------------------------------
\subsection{Tasso forward: dall'intervallo finito al limite istantaneo}
% ------------------------------------------------------------

\paragraph{Forward su intervallo finito: assenza di arbitraggio.}
Il \textbf{tasso forward} $f(t_1, t_2)$ \`e il tasso ``bloccato oggi'' per il periodo
futuro $[t_1, t_2]$. Per assenza di arbitraggio, investire fino a $t_2$ deve dare lo stesso
risultato che investire fino a $t_1$ e poi reinvestire al tasso forward da $t_1$ a $t_2$:
\[
  P(0, t_2) = P(0, t_1) \cdot e^{-f(t_1,\,t_2)\,(t_2-t_1)},
\]
da cui, prendendo il logaritmo e isolando $f$:
\begin{equation}
  f(t_1, t_2) = \frac{\ln P(0,t_1) - \ln P(0,t_2)}{t_2 - t_1}.
  \label{eq:fwd-finito}
\end{equation}
Geometricamente, $f(t_1,t_2)$ \`e il \textbf{rapporto incrementale} (con segno) di
$\ln P(0,\cdot)$ sull'intervallo $[t_1,t_2]$: la pendenza media della curva
``log-sconto'' tra le due scadenze.

\paragraph{Forward istantaneo: limite per $t_2 \to t_1$.}

\begin{definizione}[Tasso forward istantaneo]
\label{def:forward}
Facendo tendere $t_2 \to t_1 = t$ nella~\eqref{eq:fwd-finito} si ottiene il
\textbf{tasso forward istantaneo}:
\begin{equation}
  f(t) = \lim_{t_2 \to t} f(t, t_2)
       = -\frac{d}{dt}\ln P(0,t)
       = -\frac{P'(0,t)}{P(0,t)},
  \label{eq:fwd-ist}
\end{equation}
cio\`e la \textbf{derivata del log-fattore di sconto} cambiata di segno.
\end{definizione}

$f(t)\,dt$ rappresenta il tasso applicato a un prestito infinitesimale stipulato in $t$
e rimborsato in $t+dt$: il forward istantaneo \`e il ``tasso marginale'' o ``tasso locale''
della struttura a termine.

% ------------------------------------------------------------
\subsection{La relazione che unifica tutto}
% ------------------------------------------------------------

Integrando la definizione~\eqref{eq:fwd-ist} si recupera il fattore di sconto:
\begin{equation}
  P(0,t) = \exp\!\left(-\int_0^t f(s)\,ds\right).
  \label{eq:B-da-fwd}
\end{equation}
Sostituendo~\eqref{eq:B-da-fwd} in~\eqref{eq:spot-inv} si ottiene la relazione
fondamentale tra spot e forward:
\begin{equation}
  r(t) = \frac{1}{t}\int_0^t f(s)\,ds.
  \label{eq:spot-media-fwd}
\end{equation}
\begin{keybox}[Lo spot \`e la media dei forward]
L'equazione~\eqref{eq:spot-media-fwd} dice che il \textbf{tasso spot $r(t)$ \`e la media
aritmetica del forward istantaneo} sull'intervallo $[0,t]$. Una media \emph{reagisce con
ritardo} rispetto al valore puntuale: se il forward cresce, lo spot cresce pi\`u lentamente;
se il forward converge all'$\UFR$ dall'alto (come in Figura~\ref{fig:motivazione}),
lo spot ci arriver\`a \emph{dopo}. Questo spiega perch\'e in figura il forward raggiunge
l'asintoto a scadenze sensibilmente pi\`u brevi dello spot.
\end{keybox}

\medskip
La tabella seguente riassume le tre rappresentazioni equivalenti.

\begin{center}
\renewcommand{\arraystretch}{1.4}
\begin{tabular}{lll}
\toprule
\textbf{Oggetto} & \textbf{Formula} & \textbf{Lettura} \\
\midrule
Fattore di sconto & $P(0,t)$ & valore oggi di $1$\,\texteuro{} in $t$ \\
Tasso spot (continuo) & $r(t) = -\ln P(0,t)/t$ & tasso medio da $0$ a $t$ \\
Tasso forward istant. & $f(t) = -\tfrac{d}{dt}\ln P(0,t)$ & tasso marginale in $t$ \\
Ricostruzione & $P(0,t) = e^{-\int_0^t f(s)\,ds}$ & $P$ da integrale del forward \\
\bottomrule
\end{tabular}
\end{center}

\subsection{Tassi par e dati di input}

Gli strumenti effettivamente quotati sul mercato non sono zero-coupon: i prezzi
direttamente osservabili sono quelli di contratti a cedola multipla. Il concetto chiave
\`e il \textbf{tasso par}.

Considera un'obbligazione con \textbf{valore nominale} $1$, cedola annua fissa $s$ pagata
alle date $j = 1, 2, \ldots, T$, e rimborso del nominale a scadenza $T$. Il suo
\textbf{prezzo} oggi \`e la somma dei valori attuali di ogni flusso:
\[
  \text{Prezzo}(s, T)
  = s \sum_{j=1}^{T} P(0,j) + P(0,T).
\]
Un titolo che quota \emph{alla pari} ha prezzo esattamente uguale al valore nominale.

\begin{definizione}[Tasso par]
\label{def:par}
Il \textbf{tasso par} $s_T$ di scadenza $T$ \`e il tasso cedola che rende il prezzo
dell'obbligazione uguale al suo valore nominale ($= 1$):
\[
  \text{Prezzo}(s_T, T) = 1
  \;\;\Longleftrightarrow\;\;
  s_T \sum_{j=1}^{T} P(0,j) + P(0,T) = 1
  \;\;\Longrightarrow\;\;
  s_T = \frac{1 - P(0,T)}{\displaystyle\sum_{j=1}^{T} P(0,j)}.
\]
\end{definizione}

\begin{osservazione}
Un \textbf{Interest Rate Swap} (IRS) a cedola annua con scadenza $T$ pu\`o essere visto
come lo scambio della gamba variabile (floating) con la gamba fissa (cedola $s$). Per
non-arbitrage la gamba variabile vale $1$ alla data di stipula (si ``rinnova'' a ogni data
cedola), quindi il tasso fisso che rende nullo il valore dello swap all'emissione \`e
\emph{esattamente} il tasso par della Definizione~\ref{def:par}. I contratti swap IRS sono
gli strumenti liquidi quotati sul mercato a cui applicare questa formula.
\end{osservazione}

I tassi swap di mercato incorporano tuttavia un piccolo premio per il rischio di
controparte, che deve essere rimosso prima della calibrazione (cfr.\ Sez.~\ref{sec:cra}).
Il dato di input per il metodo Smith--Wilson \`e quindi il \textbf{tasso par
after-CRA}:
\begin{equation}
  s_T^{\rm EIOPA} = \frac{1 - P(0,T)}{\sum_{j=1}^{T} P(0,j)} + \CRA .
  \label{eq:par}
\end{equation}
La relazione~\eqref{eq:par} lega i dati osservati (i par rate) all'oggetto da ricostruire
(la funzione di sconto $P(0,\cdot)$), ed \`e il \emph{vincolo} che la curva Smith--Wilson
dovr\`a soddisfare ai nodi di mercato.

\begin{keybox}[Dati di input: IRS EURIBOR]
Usiamo i par rate degli \textbf{Interest Rate Swap EUR contro EURIBOR~6M} (ticker Bloomberg
\texttt{EUSA<tenor> CMPL Curncy}), file \texttt{dati/tickersEUSA\_v3.xlsx}. Le scadenze sono
le \textbf{15 scadenze DLT ufficiali} EUR SWP valide dal 30~giugno al 31~dicembre~2025:
\[
  \{1,2,3,4,5,6,7,8,9,10,11,12,13,15,20\}\ \text{anni}.
\]
\end{keybox}

La Figura~\ref{fig:spot-fwd-par} mostra, per il mese di riferimento (dicembre~2025), lo spot
e il forward ricostruiti insieme ai 15 par rate di input.

\begin{figure}[H]
\centering
\includegraphics[width=0.92\textwidth]{../output/eiopa_rfr_sw_var/fig2_spot_forward_par.pdf}
\caption{Spot $r(0,T)$, forward $f(0,T)$ e par rate \texttt{EUSA*} (punti neri) ai 15 nodi DLT,
dicembre~2025. La curva Smith--Wilson interpola i par rate e produce uno spot e un forward
$\mathcal{C}^\infty$ su tutto $[0,\CP]$.}
\label{fig:spot-fwd-par}
\end{figure}

% ============================================================
\section{Struttura matematica del metodo Smith--Wilson}
% ============================================================

L'idea di Smith--Wilson \`e cercare la funzione di sconto in una \textbf{famiglia
parametrica} costruita attorno all'$\UFR$, scegliendo i parametri in modo da
\textbf{riprezzare esattamente} gli strumenti di mercato. La struttura \`e quella di un
problema di \emph{interpolazione con kernel}, che ammette una lettura \textbf{variazionale}
(minima norma) particolarmente istruttiva.

\subsection{Il kernel di Wilson}

\begin{definizione}[Cuore e funzione di Wilson, \cite{eiopa2025} sez.~9.7]
Posto $\omega = \ln(1+\UFR)$ (intensit\`a dell'$\UFR$), si definisce il \textbf{cuore di
Wilson}
\[
  H(u,v) = \alpha\,\min(u,v) - e^{-\alpha\,\max(u,v)}\,\sinh\!\bigl(\alpha\,\min(u,v)\bigr),
\]
e la \textbf{funzione di Wilson} $W(u,v) = e^{-\omega(u+v)}\,H(u,v)$, dove $\alpha>0$ \`e un
parametro di \emph{velocit\`a di convergenza}.
\end{definizione}

La funzione di sconto viene cercata nella forma
\begin{equation}
  P(v) = e^{-\omega v}\Bigl(1 + \textstyle\sum_{j} (Qb)_j\, H(v, u_j)\Bigr),
  \label{eq:sw-P}
\end{equation}
cio\`e come \textbf{combinazione lineare} delle funzioni di Wilson centrate nei nodi di
pagamento $u_j$, pi\`u il termine di base $e^{-\omega v}$ che codifica la convergenza
all'$\UFR$. La Figura~\ref{fig:kernel} mostra le funzioni $W(\cdot,u_j)$ per alcuni nodi: sono
le ``basi'' con cui si costruisce la curva.

\begin{figure}[H]
\centering
\includegraphics[width=0.88\textwidth]{../output/eiopa_rfr_sw_var/fig3_kernel_wilson.pdf}
\caption{Funzioni di Wilson $W(t,u_j)$ centrate nei nodi $u_j \in \{1,5,10,20\}$. Ogni funzione
\`e localizzata attorno al proprio nodo e si smorza verso zero: la curva Smith--Wilson \`e una
loro combinazione lineare.}
\label{fig:kernel}
\end{figure}

\subsection{Lettura variazionale: interpolazione di minima norma}

% ----------------------------------------------------------
\paragraph{I. Spazio di Hilbert.}
% ----------------------------------------------------------

Cercare la ``curva pi\`u liscia'' che soddisfa dei vincoli di mercato \`e un problema di
\textbf{ottimizzazione in dimensione infinita}: la variabile \`e una funzione, non un
vettore finito. Per garantire l'esistenza del minimo serve uno spazio di funzioni dotato
di una struttura geometrica abbastanza ricca --- uno \emph{spazio di Hilbert}.

\begin{definizione}[Spazio di Hilbert]
\label{def:hilbert}
Uno \textbf{spazio di Hilbert} $(\mathcal{H},\langle\cdot,\cdot\rangle)$ \`e uno spazio
vettoriale (reale) dotato di un \textbf{prodotto scalare}
$\langle\cdot,\cdot\rangle\colon\mathcal{H}\times\mathcal{H}\to\mathbb{R}$ che induce la
norma $\|f\|=\sqrt{\langle f,f\rangle}$, e che \`e \textbf{completo} rispetto a tale
norma (ogni successione di Cauchy converge a un elemento di $\mathcal{H}$).
\end{definizione}

\begin{osservazione}
L'esempio canonico \`e $L^2([0,\infty))$ con
$\langle f,g\rangle = \int_0^\infty f(v)\,g(v)\,dv$. La completezza garantisce che ogni
problema di minimo su un sottospazio chiuso convesso ammetta \textbf{soluzione unica}
(teorema di proiezione). Tuttavia $L^2$ non controlla la \emph{valutazione puntuale}:
due funzioni possono coincidere in norma pur essendo diverse in ogni punto. Per
l'interpolazione occorre qualcosa di pi\`u.
\end{osservazione}

% ----------------------------------------------------------
\paragraph{II. Nucleo riproducente e RKHS.}
% ----------------------------------------------------------

Un \textbf{spazio di Hilbert a nucleo riproducente} (Reproducing Kernel Hilbert Space,
RKHS) \`e uno spazio di Hilbert di funzioni in cui la valutazione puntuale
$f \mapsto f(v)$ \`e un funzionale lineare \emph{continuo} per ogni $v$. Questa
propriet\`a \`e codificata da un \textbf{kernel}.

\begin{definizione}[Kernel simmetrico definito positivo]
\label{def:kernel-sdp}
Una funzione $K\colon X\times X\to\mathbb{R}$ \`e un \textbf{kernel simmetrico definito
positivo} (s.d.p.) se:
\begin{enumerate}
  \item $K(u,v) = K(v,u)$ per ogni $u,v\in X$ \quad (simmetria);
  \item per ogni scelta finita $v_1,\dots,v_m\in X$ e $c_1,\dots,c_m\in\mathbb{R}$
        vale $\displaystyle\sum_{i,j=1}^m c_i\,c_j\,K(v_i,v_j) \ge 0$, con uguaglianza
        solo se $c_1=\cdots=c_m=0$ \quad (definita positivit\`a stretta).
\end{enumerate}
\end{definizione}

\begin{definizione}[RKHS --- spazio di Hilbert a nucleo riproducente]
\label{def:rkhs}
Uno spazio di Hilbert $\mathcal{H}_K$ di funzioni su $X$ ha \textbf{nucleo riproducente}
$K$ se:
\begin{enumerate}
  \item $K(\cdot,v)\in\mathcal{H}_K$ per ogni $v\in X$;
  \item $\langle f,\,K(\cdot,v)\rangle_{\mathcal{H}_K} = f(v)$
        per ogni $f\in\mathcal{H}_K$ e ogni $v\in X$
        \quad(\textbf{propriet\`a riproducente}).
\end{enumerate}
\end{definizione}

La propriet\`a riproducente ha una conseguenza immediata: per la disuguaglianza di
Cauchy--Schwarz,
\[
  |f(v)| = |\langle f,\,K(\cdot,v)\rangle_{\mathcal{H}_K}|
  \;\le\; \|f\|_{\mathcal{H}_K}\cdot\|K(\cdot,v)\|_{\mathcal{H}_K},
\]
quindi la valutazione puntuale \`e continua in norma. Questo rende l'interpolazione
(imporre $f(v_k) = y_k$) un'operazione ben posta nello spazio.

\begin{teorema}[Moore--Aronszajn, \cite{aronszajn1950}]
\label{thm:rkhs}
C'\`e una corrispondenza biunivoca tra kernel s.d.p.\ e RKHS: a ogni kernel s.d.p.\ $K$
su $X$ corrisponde un \textbf{unico} RKHS $\mathcal{H}_K$ con nucleo riproducente $K$, e
viceversa ogni RKHS ha un unico nucleo riproducente.
\end{teorema}

% ----------------------------------------------------------
\paragraph{III. Collegamento con Smith--Wilson.}
% ----------------------------------------------------------

Il cuore di Wilson $H$ \`e un kernel s.d.p.\ su $[0,\infty)$: per ogni scelta finita di
punti la matrice $[H(v_i,v_j)]$ \`e definita positiva. Per il
Teorema~\ref{thm:rkhs} genera quindi un unico RKHS $\mathcal{H}_H$, la cui norma ha la
forma esplicita
\[
  \|f\|_{\mathcal{H}_H}^2
  = \int_0^\infty \bigl(f''(v) + \omega^2\,f(v)\bigr)^2\,dv.
\]
Questa norma penalizza le funzioni che si discostano dall'asintotico $e^{-\omega v}$
(convergenza all'$\UFR$): \`e la misura di ``ruvidezza'' scelta da EIOPA.

I vincoli di riprezzamento~\eqref{eq:par} definiscono, per ogni scadenza $k$, il
funzionale lineare $\ell_k\colon P \mapsto s_k^{\rm EIOPA}(P)$. Per la propriet\`a
riproducente ogni $\ell_k$ \`e continuo in $\mathcal{H}_H$. Il metodo Smith--Wilson
risolve il problema variazionale:
\begin{equation}
  P^* = \operatorname*{arg\,min}_{P\,\in\,\mathcal{H}_H} \|P\|_{\mathcal{H}_H}
  \qquad \text{s.t.}\quad
  \ell_k(P) = \hat{s}_k,\quad k = 1,\dots,n.
  \label{eq:varmin}
\end{equation}

\begin{teorema}[Esistenza e unicità, \cite{berlinet2004} cap.~2]
\label{thm:existunique}
Il problema~\eqref{eq:varmin} ammette \textbf{soluzione unica}. L'esistenza segue dalla
completezza di $\mathcal{H}_H$ (spazio di Hilbert, Def.~\ref{def:hilbert}); l'unicit\`a
dalla \textbf{stretta convessit\`a} della norma e dalla non-degenerazione dei vincoli,
garantita dalla definita positivit\`a stretta di $H$ sui nodi di calibrazione.
\end{teorema}

\begin{teorema}[Representer theorem, \cite{wahba1990}, \cite{berlinet2004}]
\label{thm:representer}
La soluzione $P^*$ di~\eqref{eq:varmin} appartiene allo spazio di dimensione finita
\[
  \mathcal{V}_n = \operatorname{span}\bigl\{H(\cdot,\,u_j)\bigr\}_{j=1}^{n},
\]
dove $u_1,\dots,u_n$ sono le date di pagamento degli strumenti. Il problema
infinito-dimensionale si riduce quindi a determinare $n$ coefficienti reali, cio\`e
esattamente il sistema lineare~\eqref{eq:sistema}.
\end{teorema}

\begin{proposizione}[Caratterizzazione di minima norma]
\label{prop:minnorm}
Tra tutte le funzioni di sconto che soddisfano~\eqref{eq:par} ai nodi osservati, la
soluzione Smith--Wilson \`e l'unica con \textbf{norma $\|\cdot\|_{\mathcal{H}_H}$ minima}.
\end{proposizione}

\begin{keybox}[Il filo logico: da Hilbert al sistema $n\times n$]
\begin{enumerate}
  \item \textbf{Spazio di Hilbert $\mathcal{H}_H$}: il kernel di Wilson $H$ (s.d.p.) genera
        per Moore--Aronszajn uno spazio completo con prodotto scalare. La norma misura
        la ``ruvidezza'' e forza la convergenza all'$\UFR$.
  \item \textbf{Ottimizzazione in $\mathcal{H}_H$}: il problema ``curva pi\`u liscia che
        passa per i dati'' \`e convesso vincolato in dimensione infinita; per il
        Teorema~\ref{thm:existunique} ha soluzione unica.
  \item \textbf{Representer theorem}: la soluzione vive in $\mathcal{V}_n$ (dim.\ $n=15$).
        Il problema infinito collassa al sistema lineare~\eqref{eq:sistema} --- lo stesso
        principio delle \emph{spline} cubiche, ma con la norma di Wilson.
\end{enumerate}
\end{keybox}

\subsection{Lo stesso risultato con la sola algebra lineare: minima energia e proiezione}

La lettura variazionale precedente richiede l'analisi funzionale (spazi di Hilbert,
nuclei riproducenti). Non \`e per\`o necessaria: lo \textbf{stesso sistema lineare} di
calibrazione si ottiene restando in \textbf{dimensione finita}, con i soli strumenti
dell'algebra lineare --- matrici simmetriche definite positive (SPD), forme quadratiche,
prodotti scalari indotti da una matrice SPD, moltiplicatori di Lagrange e proiezioni
ortogonali. Procediamo in cinque passi.

% ----------------------------------------------------------
\paragraph{Passo 1 --- Impostazione: l'incognita \`e un vettore finito.}
% ----------------------------------------------------------
Cerchiamo la funzione di sconto nella forma~\eqref{eq:sw-P}, cio\`e
\[
  P(v) = e^{-\omega v}\Bigl(1 + \sum_{j=1}^{m}\zeta_j\,H(v,u_j)\Bigr),
\]
dove $u_1,\dots,u_m$ sono le $m=\LLP$ date di pagamento e $n$ \`e il numero di strumenti di
input. Una volta fissato il kernel $H$, la curva \`e individuata \emph{soltanto} dal
\textbf{vettore finito} di coefficienti $\zeta = (\zeta_1,\dots,\zeta_m)\transp \in
\mathbb{R}^m$: il problema diventa allora finito-dimensionale, ambientato nel sottospazio
$\mathcal{V}_n = \operatorname{span}\{H(\cdot,u_j)\}_{j=1}^{m}$. Tutto ci\`o che segue \`e
algebra lineare su $\mathbb{R}^m$.

% ----------------------------------------------------------
\paragraph{Passo 2 --- Una matrice SPD definisce un prodotto scalare: l'energia.}
% ----------------------------------------------------------
Il punto di partenza \`e un fatto generale di algebra lineare, che enunciamo con
precisione perch\'e \`e il cardine di tutta la costruzione.

\begin{osservazione}[Una matrice SPD induce un prodotto scalare]
\label{oss:spd-prodotto}
Sia $\mathbf{A}\in\mathbb{R}^{m\times m}$ \textbf{simmetrica definita positiva} (SPD).
Allora la forma $\langle x,y\rangle_{\mathbf{A}} := x\transp \mathbf{A}\, y$ \`e un
\textbf{prodotto scalare} su $\mathbb{R}^m$, perch\'e ne soddisfa i tre assiomi:
\begin{enumerate}
  \item \textbf{simmetria}: $\langle x,y\rangle_{\mathbf{A}} = \langle y,x\rangle_{\mathbf{A}}$,
        conseguenza di $\mathbf{A} = \mathbf{A}\transp$ (\,$\mathbf{A}$ \emph{simmetrica}\,);
  \item \textbf{bilinearit\`a}: $\langle \cdot,\cdot\rangle_{\mathbf{A}}$ \`e lineare in
        ciascun argomento, perch\'e tale \`e il prodotto matrice--vettore;
  \item \textbf{definita positivit\`a}: $\langle x,x\rangle_{\mathbf{A}} = x\transp\mathbf{A}\,x > 0$
        per ogni $x\neq 0$, conseguenza di $\mathbf{A}$ \emph{definita positiva}.
\end{enumerate}
Di conseguenza $\|x\|_{\mathbf{A}} := \sqrt{x\transp \mathbf{A}\, x}$ \`e una \textbf{norma}.
\end{osservazione}

Applichiamo l'Osservazione~\ref{oss:spd-prodotto} alla \textbf{matrice di Wilson}
$\Hmat = [H(u_i,u_j)]_{i,j=1}^{m}$: essa \`e la \textbf{matrice di Gram} delle sezioni del
kernel $H(\cdot,u_j)$ e, poich\'e $H$ \`e un kernel simmetrico definito positivo
(Def.~\ref{def:kernel-sdp}), \`e essa stessa \textbf{SPD}. Per
l'Osservazione~\ref{oss:spd-prodotto}, $\Hmat$ induce su $\mathbb{R}^m$ il prodotto scalare
$\langle x,y\rangle_{\Hmat} = x\transp\Hmat\,y$. Chiamiamo allora \textbf{energia}
(ruvidezza) della curva la forma quadratica
\[
  E(\zeta) = \zeta\transp \Hmat\, \zeta = \|\zeta\|_{\Hmat}^2 \ \ge 0 .
\]
\`E la stessa quantit\`a della norma $\|\cdot\|_{\mathcal{H}_H}$ della
Proposizione~\ref{prop:minnorm}, ristretta a $\mathcal{V}_n$, ma qui \`e calcolata con la
\textbf{sola matrice di Gram}: nessun integrale, nessuno spazio infinito-dimensionale.

% ----------------------------------------------------------
\paragraph{Passo 3 --- I vincoli di mercato sono affini in $\zeta$.}
% ----------------------------------------------------------
Valutando~\eqref{eq:sw-P} ai nodi si ottiene il vettore dei fattori di sconto
$\mathbf{P} = \bigl(P(u_1),\dots,P(u_m)\bigr)\transp$, cio\`e, con $d_j = e^{-\omega u_j}$,
\[
  \mathbf{P} = d + \operatorname{diag}(d)\,\Hmat\,\zeta .
\]
Imporre il riprezzamento esatto degli strumenti, $\mathbf{C}\transp \mathbf{P} = p$, e posto
$\Qmat = \operatorname{diag}(d)\,\mathbf{C}$ e $q = \mathbf{C}\transp d$, d\`a la condizione
\[
  \Qmat\transp \Hmat\, \zeta = p - q .
\]
Si tratta di un \textbf{sistema lineare} di $n$ equazioni nelle $m$ incognite $\zeta$:
i suoi vettori soluzione formano un \textbf{insieme affine} (non vuoto) di $\mathbb{R}^m$.

% ----------------------------------------------------------
\paragraph{Passo 4 --- Il problema: minima energia su un affine.}
% ----------------------------------------------------------
Tra tutte le curve che riprezzano gli strumenti scegliamo quella di energia minima:
\begin{equation}
  \min_{\zeta\,\in\,\mathbb{R}^m} \ \zeta\transp \Hmat\, \zeta
  \qquad \text{s.t.}\quad
  \Qmat\transp \Hmat\, \zeta = p - q .
  \label{eq:minenergia}
\end{equation}
Questo problema \`e \textbf{ben posto} gi\`a prima di calcolarne la soluzione: la funzione
obiettivo $\zeta\mapsto\zeta\transp\Hmat\zeta$ \`e una forma quadratica
\textbf{strettamente convessa} (perch\'e $\Hmat$ \`e definita positiva) ristretta a un
insieme affine non vuoto. Un tale problema ammette uno e un solo punto di minimo.

\begin{proposizione}[Minima norma, versione finito-dimensionale]
\label{prop:minnorm-finita}
Il problema~\eqref{eq:minenergia} ammette \textbf{soluzione unica}. L'esistenza e
l'unicit\`a seguono dalla sola \textbf{definita positivit\`a} di $\Hmat$ (stretta
convessit\`a dell'energia) su un vincolo affine: nessuna nozione di analisi funzionale \`e
necessaria.
\end{proposizione}

% ----------------------------------------------------------
\paragraph{Passo 5 --- Soluzione (Lagrange) e lettura geometrica (proiezione).}
% ----------------------------------------------------------
Per calcolare il minimo introduciamo i moltiplicatori di Lagrange $b\in\mathbb{R}^n$, uno
per ciascuno degli $n$ vincoli. La condizione di stazionariet\`a
$\nabla_\zeta\bigl(\zeta\transp\Hmat\zeta - 2\,b\transp(\Qmat\transp\Hmat\zeta - (p-q))\bigr)=0$
d\`a $\Hmat\zeta = \Hmat\Qmat\,b$, da cui (essendo $\Hmat$ invertibile)
\[
  \zeta = \Qmat\, b .
\]
Sostituendo nel vincolo del Passo~3 si ritrova \textbf{esattamente} il sistema di
calibrazione~\eqref{eq:sistema}:
\[
  \bigl(\Qmat\transp \Hmat\, \Qmat\bigr)\, b = p - q .
\]
I coefficienti $b$ del metodo Smith--Wilson non sono dunque altro che i
\textbf{moltiplicatori di Lagrange} del problema di minima energia~\eqref{eq:minenergia}.

La relazione $\zeta = \Qmat\,b$ ha un significato \textbf{geometrico}: il vettore ottimo dei
coefficienti appartiene a $\operatorname{range}(\Qmat)$, il \textbf{sottospazio generato
dalle colonne di} $\Qmat$ (le sezioni del kernel pesate dai flussi di cassa). La soluzione
\`e quindi la \textbf{proiezione $\Hmat$-ortogonale} della curva base $e^{-\omega v}$
sull'insieme delle curve che riprezzano gli strumenti: tra tutte le correzioni ammissibili
si sceglie quella di lunghezza minima nella metrica indotta da $\Hmat$. \`E l'analogo
finito-dimensionale, e completamente elementare, del principio di minima norma della
Proposizione~\ref{prop:minnorm}.

\begin{keybox}[Smith--Wilson senza analisi funzionale]
La stessa curva si ottiene con la sola algebra lineare:
\begin{enumerate}
  \item la matrice di Wilson $\Hmat$ \`e \textbf{SPD} $\Rightarrow$ definisce un prodotto
        scalare e una nozione di \emph{energia} $\zeta\transp\Hmat\zeta$;
  \item si minimizza l'energia sotto i \textbf{vincoli lineari} di repricing
        (\,\eqref{eq:minenergia}\,);
  \item i \textbf{moltiplicatori di Lagrange} $b$ risolvono il sistema SPD
        $(\Qmat\transp\Hmat\Qmat)b = p-q$, di ordine $n$ (Cholesky);
  \item geometricamente, la curva \`e la \textbf{proiezione ortogonale} (nella metrica
        $\Hmat$) sul sottospazio generato dalle funzioni $H$.
\end{enumerate}
\`E la lettura adatta a chi conosce l'algebra lineare ma non ancora l'analisi funzionale:
cfr.\ \cite{boyd2018} per i minimi quadrati vincolati, \cite{fasshauer2007} per
l'interpolazione con kernel definiti positivi, \cite{lageras2016} per il caso Smith--Wilson.
\end{keybox}

\subsection{Il sistema lineare di calibrazione}

Imporre il riprezzamento esatto degli strumenti conduce, dopo aver scritto i flussi di cassa
nella matrice $\mathbf{C}$ ($m\times n$, con $m=\LLP$ date di pagamento e $n$ strumenti) e
posto $d_j = e^{-\omega u_j}$, $\Qmat = \mathrm{diag}(d)\,\mathbf{C}$,
$\Hmat = [H(u_i,u_j)]$, $q = \mathbf{C}\transp d$, al sistema lineare
(\cite{eiopa2025} sez.~9.8--9.10)
\begin{equation}
  \bigl(\Qmat\transp \Hmat\, \Qmat\bigr)\, b = p - q,
  \qquad p = \mathbf{1}\ \text{(par swap)} .
  \label{eq:sistema}
\end{equation}
Risolto per $b$, i coefficienti $Qb = \Qmat b$ entrano direttamente
nella~\eqref{eq:sw-P}. La matrice del sistema ha una struttura che ne garantisce la buona
risolubilit\`a numerica.

\begin{teorema}[Sistema SPD e fattorizzazione di Cholesky]
\label{thm:spd}
La matrice $\Qmat\transp \Hmat\, \Qmat \in \mathbb{R}^{n\times n}$ \`e \textbf{simmetrica
definita positiva}. Di conseguenza il sistema~\eqref{eq:sistema} ammette \textbf{soluzione
unica}, calcolabile in modo stabile tramite la \textbf{fattorizzazione di Cholesky}
$\Qmat\transp \Hmat\,\Qmat = R\transp R$ con $R$ triangolare superiore.
\end{teorema}

La Figura~\ref{fig:QHQ} visualizza la matrice $\Qmat\transp\Hmat\Qmat$ per il mese di
riferimento: simmetrica e con autovalore minimo strettamente positivo.

\begin{figure}[H]
\centering
\includegraphics[width=0.72\textwidth]{../output/eiopa_rfr_sw_var/fig4_matrice_QHQ.pdf}
\caption{Matrice del sistema $\Qmat\transp\Hmat\Qmat$ ($n=15$), dicembre~2025. \`E simmetrica
e definita positiva ($\lambda_{\min} > 0$): la fattorizzazione di Cholesky \`e applicabile ed
\`e circa due volte pi\`u efficiente di una LU generica.}
\label{fig:QHQ}
\end{figure}

\begin{teorema}[Rappresentazione Smith--Wilson]
\label{thm:rappresentazione}
Risolto il sistema~\eqref{eq:sistema}, la funzione di sconto~\eqref{eq:sw-P} riprezza
esattamente tutti gli strumenti di input ed \`e di classe $\mathcal{C}^\infty$ su $[0,\CP]$.
Da essa si ricavano spot e forward tramite
\[
  r(0,v) = P(v)^{-1/v}-1, \qquad
  f(0,v) = \omega - \frac{\sum_j (Qb)_j\,\partial_v H(v,u_j)}{1 + \sum_j (Qb)_j\, H(v,u_j)} .
\]
\end{teorema}

\subsection{Calibrazione di $\alpha$: un problema di ricerca di zeri}

Il parametro $\alpha$ governa la velocit\`a di convergenza all'$\UFR$. EIOPA lo fissa con un
criterio preciso (\cite{eiopa2025} sez.~9.14): si sceglie il \textbf{pi\`u piccolo}
$\alpha \ge \alpha_{\min}=0.05$ tale che il forward al Convergence Point disti dall'$\UFR$
meno di una tolleranza $\tau = 1$ bp:
\begin{equation}
  \alpha^\star = \min\Bigl\{\alpha \ge \alpha_{\min} : \bigl|f(0,\CP;\alpha) - \omega\bigr|
  \le \tau \Bigr\}.
  \label{eq:criterio-alpha}
\end{equation}
Definita $g(\alpha) = |f(0,\CP;\alpha) - \omega| - \tau$, la~\eqref{eq:criterio-alpha} \`e un
\textbf{problema di ricerca di zeri}, risolvibile con il metodo di \textbf{bisezione} (robusto)
o di \textbf{Newton} (rapido). La Figura~\ref{fig:alpha} mostra $g$ per il mese di riferimento.

\begin{figure}[H]
\centering
\includegraphics[width=0.9\textwidth]{../output/eiopa_rfr_sw_var/fig5_calibrazione_alpha.pdf}
\caption{Calibrazione di $\alpha$ come ricerca di zeri (dicembre~2025). La curva \`e
$f(0,\CP;\alpha)-\omega$ in bps; la banda grigia \`e la tolleranza $\pm 1$ bp. Il valore
$\alpha^\star \approx 0.064$ \`e il pi\`u piccolo $\alpha \ge \alpha_{\min}$ che entra nella
banda.}
\label{fig:alpha}
\end{figure}

\begin{teorema}[Convergenza all'UFR]
\label{thm:convergenza}
Per ogni $\alpha>0$ il forward Smith--Wilson converge all'intensit\`a dell'$\UFR$:
\[
  \lim_{v\to\infty} f(0,v) = \omega = \ln(1+\UFR).
\]
\end{teorema}

\begin{algobox}[Costruzione della curva Smith--Wilson (in R)]
\begin{Verbatim}[fontsize=\small]
# INPUT: par rate EUSA after-CRA r_k ai nodi T_mkt = {1..13,15,20}

# 1. Matrice dei flussi C (m x n), m = LLP date di pagamento annuali
C <- build_C(T_mkt, r)                 # C[j,k]=r_k (j<T_k), C[T_k,k]=1+r_k

# 2. Oggetti del sistema (dato alpha)
d   <- exp(-omega * (1:LLP)); Q <- C * d
H   <- outer(1:LLP, 1:LLP, H_heart, a = alpha)
QHQ <- t(Q) %*% H %*% Q                # SIMMETRICA DEFINITA POSITIVA

# 3. Soluzione via Cholesky (QHQ = R'R)
R  <- chol(QHQ)
b  <- backsolve(R, forwardsolve(t(R), p - q))
Qb <- as.numeric(Q %*% b)

# 4. alpha* via ricerca di zeri: |f(CP,alpha) - omega| <= 1bp, alpha >= 0.05
alpha_star <- alpha_star_for_mat(r, T_mkt)

# 5. Curva: P(v) = e^{-omega v}(1 + sum_j Qb_j H(v,u_j))
spot <- sw_spot_ann(v, Qb, alpha_star)
\end{Verbatim}
\end{algobox}

% ============================================================
\section{Ricostruzione della curva: giugno--dicembre 2025}
\label{sec:ricostruzione}
% ============================================================

Mettiamo alla prova il metodo ricostruendo la curva EIOPA EUR sui \textbf{sette mesi} da
giugno a dicembre~2025, la finestra in cui il set di scadenze $\{1,\dots,13,15,20\}$ \`e quello
ufficiale. Per ogni mese:
\begin{enumerate}
  \item leggiamo i par rate \texttt{EUSA*} e sottraiamo il $\CRA$;
  \item calibriamo Smith--Wilson (Cholesky + ricerca di $\alpha$);
  \item confrontiamo la curva spot ricostruita con quella \textbf{ufficiale} pubblicata da
        EIOPA (file \texttt{dati/eiopa\_zips/}).
\end{enumerate}

\subsection{Confronto grafico}

\begin{figure}[H]
\centering
\includegraphics[width=0.92\textwidth]{../output/eiopa_rfr_sw_var/fig6_overlay_ricostruzione.pdf}
\caption{Curva spot ricostruita con Smith--Wilson su \texttt{EUSA*} (tratteggiata) vs curva
ufficiale EIOPA (continua), per i sette fine-mese da giugno a dicembre~2025. Le due curve sono
visivamente sovrapposte su tutta la struttura a termine.}
\label{fig:overlay}
\end{figure}

La Figura~\ref{fig:banda} sintetizza lo scarto: l'inviluppo min--max dei residui sui sette mesi
resta entro pochi bps in zona liquida e si mantiene contenuto anche in estrapolazione.

\begin{figure}[H]
\centering
\includegraphics[width=0.92\textwidth]{../output/eiopa_rfr_sw_var/fig7_banda_residui.pdf}
\caption{Banda delle differenze (ricostruita $-$ ufficiale) sui sette mesi: regione
ombreggiata = inviluppo min--max per scadenza, linea = media. In zona liquida lo scarto
massimo \`e di poco superiore al bp.}
\label{fig:banda}
\end{figure}

\begin{figure}[H]
\centering
\includegraphics[width=0.92\textwidth]{../output/eiopa_rfr_sw_var/fig8_barre_per_tenor.pdf}
\caption{Differenza ricostruita $-$ ufficiale (bps) ai 15 nodi DLT, per ciascun mese. Barre
arancioni: scarto positivo; barre blu: negativo. Quasi tutte le barre stanno entro $\pm 2$ bps.}
\label{fig:barre}
\end{figure}

\subsection{Risultati quantitativi}

La Tabella~\ref{tab:rmse} riporta, per ogni mese, l'$\alpha$ di criterio e quello ufficiale e
gli RMSE dello scarto in zona liquida ($T\le\LLP$) e in estrapolazione ($T>\LLP$).

\begin{table}[H]
\centering\small
\caption{Ricostruzione Smith--Wilson da \texttt{EUSA*} vs curva ufficiale EIOPA, giugno--dicembre~2025.
Prodotta da \texttt{R/eiopa\_rfr\_sw\_variazionale.R}, salvata in
\texttt{output/eiopa\_rfr\_sw\_var/tab\_ricostruzione\_var.csv}.}
\label{tab:rmse}
\renewcommand{\arraystretch}{1.2}
\begin{tabular}{lccccc}
\toprule
\textbf{Data} & $\alpha^\star_{\rm crit}$ & $\alpha_{\rm uff}$ &
\textbf{RMSE}$_{\le 20}$ & \textbf{RMSE}$_{>20}$ & $\max|\Delta|$\\
 & & & (bps) & (bps) & (bps)\\
\midrule
\rowcolor{blue!4}2025-06-30 & 0.0946 & 0.0936 & 1.04 & 0.38 & 1.78\\
2025-07-31 & 0.0897 & 0.0908 & 0.39 & 0.16 & 0.73\\
\rowcolor{blue!4}2025-08-29 & 0.0752 & 0.0781 & 0.27 & 0.51 & 0.65\\
2025-09-30 & 0.0810 & 0.0822 & 0.11 & 0.36 & 0.49\\
\rowcolor{blue!4}2025-10-31 & 0.0754 & 0.0790 & 0.25 & 0.79 & 1.00\\
2025-11-28 & 0.0500 & 0.0500 & 0.40 & 1.87 & 2.20\\
\rowcolor{blue!4}2025-12-31 & 0.0635 & 0.0736 & 0.46 & 1.41 & 1.72\\
\bottomrule
\end{tabular}
\end{table}

\begin{codebox}[Ricostruzione mensile (estratto)]
\begin{Verbatim}[fontsize=\small]
for (i in seq_len(nrow(eusa))) {           # eusa = 7 mesi giu-dic 2025
  s_i  <- as.numeric(eusa[i, -1])          # par EUSA (decimali)
  off  <- read_eiopa_official(eusa$date[i])# curva ufficiale dal relativo zip
  r_i  <- s_i - off$CRA/1e4                 # after-CRA
  a    <- off$alpha                         # alpha ufficiale del mese
  Qb   <- sw_calibrate(T_mkt, r_i, a)$Qb    # Cholesky
  sp   <- sw_spot_ann(1:80, Qb, a) * 100    # spot ricostruito
  res  <- (sp - sp_ufficiale) * 100         # residuo in bps
}
\end{Verbatim}
\end{codebox}

\begin{osservazione}
In tutti i sette mesi l'$\alpha$ di criterio e quello ufficiale quasi coincidono e gli RMSE
restano di \emph{pochi} bps, sia in zona liquida sia in estrapolazione. Ci\`o conferma due
fatti: il metodo Smith--Wilson \`e numericamente esatto nel riprezzamento, e l'\textbf{IRS
EURIBOR} (\texttt{EUSA*}) \`e lo strumento di input coerente con la curva ufficiale EIOPA.
\end{osservazione}

% ============================================================
\section{Conclusioni}
% ============================================================

\begin{itemize}
  \item La curva EIOPA risk-free risolve un problema di \textbf{interpolazione ed
        estrapolazione}: seguire i tassi di mercato fino al $\LLP$ e convergere all'$\UFR$ oltre.
  \item Il metodo \textbf{Smith--Wilson} costruisce la funzione di sconto come combinazione di
        \textbf{funzioni di Wilson}; ha una lettura \textbf{variazionale} (interpolante di
        minima norma nell'RKHS del kernel), analoga a quella delle spline.
  \item Il nucleo numerico \`e un \textbf{sistema lineare SPD} risolto con \textbf{Cholesky}
        (Teorema~\ref{thm:spd}) e una \textbf{ricerca di zeri} per il parametro $\alpha$
        (criterio sez.~9.14).
  \item La ricostruzione sui sette mesi giugno--dicembre~2025 riproduce la curva ufficiale EIOPA
        entro pochi bps, validando sia l'implementazione sia la scelta dell'input di mercato.
\end{itemize}

% ============================================================
\renewcommand{\refname}{Riferimenti bibliografici}
\begin{thebibliography}{9}

\bibitem{eiopa2025}
  \textbf{EIOPA} (2025).
  \textit{Technical documentation of the methodology to derive EIOPA's risk-free interest rate
  term structures}, EIOPA-BoS-25-599.
  Sezione~9 (Extrapolation and interpolation) per il metodo Smith--Wilson.

\bibitem{solvency2}
  \textbf{Direttiva 2009/138/CE} (Solvency~II) e
  \textbf{Regolamento Delegato (UE) 2015/35}.
  Quadro normativo per la valutazione market-consistent delle riserve tecniche.

\bibitem{smithwilson2001}
  \textbf{Smith, A., Wilson, T.} (2001).
  \textit{Fitting yield curves with long term constraints}.
  Bacon \& Woodrow Research Notes.

\bibitem{quarteroni2017}
  \textbf{Quarteroni, A., Saleri, F., Gervasio, P.} (2017).
  \textit{Calcolo Scientifico}, 6a ed. Springer.
  Fattorizzazione di Cholesky (cap.~5) e metodi per la ricerca di zeri (cap.~2).

\bibitem{burden2016}
  \textbf{Burden, R.L., Faires, J.D.} (2016).
  \textit{Numerical Analysis}, 10th ed. Cengage.
  Bisezione, Newton e sistemi lineari definiti positivi.

\bibitem{boyd2018}
  \textbf{Boyd, S., Vandenberghe, L.} (2018).
  \textit{Introduction to Applied Linear Algebra: Vectors, Matrices, and Least Squares}.
  Cambridge University Press.
  Cap.~16: minimi quadrati vincolati e soluzione di minima norma via moltiplicatori di
  Lagrange. Testo liberamente disponibile (Stanford/UCLA).

\bibitem{fasshauer2007}
  \textbf{Fasshauer, G.E.} (2007).
  \textit{Meshfree Approximation Methods with MATLAB}. World Scientific.
  Interpolazione con kernel definiti positivi: l'interpolante \`e combinazione lineare delle
  sezioni del kernel e i coefficienti risolvono il sistema con la matrice di Gram SPD.

\bibitem{lageras2016}
  \textbf{Lager\aa s, A., Lindholm, M.} (2016).
  Issues with the Smith--Wilson method.
  \textit{Insurance: Mathematics and Economics}, 68, 32--41.
  Trattazione del metodo Smith--Wilson come interpolazione con kernel.

\bibitem{aronszajn1950}
  \textbf{Aronszajn, N.} (1950).
  Theory of reproducing kernels.
  \textit{Transactions of the American Mathematical Society}, 68, 337--404.

\bibitem{wahba1990}
  \textbf{Wahba, G.} (1990).
  \textit{Spline Models for Observational Data}. SIAM, Philadelphia.
  Cap.~1--2: spline di interpolazione, minima norma, Representer theorem.

\bibitem{berlinet2004}
  \textbf{Berlinet, A., Thomas-Agnan, C.} (2004).
  \textit{Reproducing Kernel Hilbert Spaces in Probability and Statistics}.
  Kluwer Academic Publishers.
  Cap.~2: esistenza e unicità dell'interpolante di minima norma nell'RKHS.

\end{thebibliography}

\end{document}
